\scanpage{347}
\mspage{86}
% archive sheet number: 346
L'idée donc, c'est d'être dans les dispositions d'un voyageur qui, tenu à un plan de voyage par étapes décrit par $\tau$, voudrait arriver aussi vite que possible à $1_{\tau'}$, en partant de $0_{\tau'}$. Il va dès la première étape, \deleted{\unsure{$\Delta$}} $\Delta_\varepsilon\mathrel{\unsure{\subset|\in}}I_1$, \inserted{se} porter jusqu'au bout de $I'_1$, regrettant ne pas pouvoir aller plus loin (à cause du plan de voyage). Il doit rester inactif jusqu'à ce qu'il soit arrivé au \deleted{\unsure{temps}} moment marqué par l'extrémité finale de $I_1$ (si celle-ci est au delà de celle de sa première flèche de transition, i.e. si $\ell_1>1$), puis il fait son deuxième grand pas jusqu'à l'extrémité terminale de $I'_2$ etc. Il peut espérer, à chaque moment de son voyage, être en avance \inserted{(ou au moins pas en retard)} sur les autres voyageurs, qui auraient suivi \unsure{des} \deleted{\unsure{$\beta$}} tactiques de progression différentes.

Si \deleted{\unsure{et}} cet espoir était justifié, cela signifierait que $F_0$ serait un plus \deleted{petit} \inserted{grand} élément de $\underline{\mathrm{Hom}}(\tau,\tau')$ (\uline{NB} toutes ces catégories $\underline{\mathrm{Hom}}$ \inserted{entre intervalles} sont ``ordonnées'' i.e. il y a au plus une flèche entre deux objets), i.e. un objet \deleted{\unsure{initial}} final. Cela impliquerait \unsure{de|que} $\underline{\mathrm{Hom}}(\tau,\tau')$ serait même \unsure{contractible|contractile} -- chose peut-être vraie, mais sûrement pas de façon aussi simpliste je crains. Je me suis d'ailleurs \unsure{coulé}~:
$F_0$ est \deleted{plus petit} sans doute bien un ``plus grand'' \deleted{\unsure{Essayer|Enfin}, \unsure{quand même, de prouver que pour toute|quand même, de prouver que pour tout} $F$, \illegible{1}} élément, mais pour l'ordre défini par la numérotation des sommets de $I_{\tau'}$ (donc $F\leq G$ sss $F(i)\leq G(i)\quad\forall i\in I_\tau$), et nullement pour l'ordre défini par la structure de catégorie de $\underline{\mathrm{Hom}}(\tau,\tau')$.
